\section{A uniform energy bound}\label{sec:energy}
The purpose of this Section is to replace the two many-variable
majorants by one-variable suprema and explicit quadratic terms. The
bound must be uniform before taking those suprema. Matched damping
and retained endpoint corrections provide this uniformity.

We retain the constants $\alpha,\beta,\gamma,\eta$, the two cases
$\kappa=2,1$, and the nonnegative majorants $\mathcal I_{\kappa,n}$ of
Equation~\eqref{eq:real-integrand-majorants}.  Put
\begin{align}
 W_\kappa(x)&=(\alpha+2\gamma)\log|x|+2\eta\log(1-x)
 -(\kappa/2+\alpha+\eta+\gamma)\log(1+x^2),\notag\\
 W_s(s)&=\beta\log s+\gamma\log(1-s),\notag\\
 D(x)&=2\gamma-\frac{2x^2}{1+x^2}.
 \label{eq:energy-fields}
\end{align}
Here $D(x)$ is a scalar function, distinct from the row polynomials $D_r(t)$.
For a real sequence $u=(u_k)_{k\ge1}$ with $|u_k|\le K r^k$ for some
$K<\infty$ and $0<r<1$, define
\begin{equation}\label{eq:energy-trial-notation}
 \|u\|_*^2=\sum_{k\ge1}\frac{u_k^2}{k},\qquad
 T(u,x)=\sum_{k\ge1}\frac{u_kT_k(x)}{k},\qquad
 S(u,x)=\sum_{k\ge1}\frac{u_kx^k}{k}.
\end{equation}
These series converge uniformly for $-1\le x\le1$.

The sequences $p,v$ below are freely chosen trial coefficients. The
inequality $-a^2\le b^2-2ab$ lets them replace two negative quadratic
sums in the logarithmic interactions by affine upper bounds, leaving
the two one-variable suprema. Each admissible choice gives a valid
bound; the certificate will supply one successful choice for each case.

\begin{proposition}\label{prop:energy-reduction}
Let $p,v$ be two such sequences, and take $\lambda=0$ in case $\kappa=2$
or any fixed $\lambda\ge0$ in case $\kappa=1$.  Uniformly over configurations
in the indicated case,
\begin{align}
 &\limsup_{N\to\infty}\ \sup_{\text{case }\kappa}
 \left(\frac{\log\mathcal I_{\kappa,n}}{n^2}
                  -\frac12\log2\right)\notag\\
 &\qquad\le (-1+\alpha+\gamma)\log2
       +\kappa\|p\|_*^2+\frac12\|v\|_*^2\notag\\
 &\qquad\quad+\sup_{\substack{-1\le x<1\\x\ne0}}
       \{W_\kappa(x)+\lambda D(x)-2\kappa T(p,x)-S(v,x)\}\notag\\
 &\qquad\quad+\sup_{0<s<1}\{W_s(s)+2T(v,s)\}.
 \label{eq:energy-dual}
\end{align}
Consequently the maximum of the right sides for the two cases bounds
\[
 \limsup_{N\to\infty}\left(n^{-2}\log|\Delta_N|-\tfrac12\log2\right).
\]
We use $\log0=-\infty$.
\end{proposition}

\begin{proof}
For a one-variable function $F$, write
$\langle F(x)\rangle=n^{-1}\sum_iF(x_i)$, and similarly for $s$.
In a double average $\langle K(x,x')\rangle_{\ne}$ we sum over ordered
pairs $i\ne j$ and divide by $n^2$; a double average without the subscript
includes every pair.  The substitution $t=2x/(1+x^2)$ gives
\begin{equation}\label{eq:energy-transform-identities}
 |t-t'|=\frac{2|x-x'|(1-xx')}{(1+x^2)(1+x'^2)},\qquad
 1-ts=\frac{1-2xs+x^2}{1+x^2},\qquad
 1-t=\frac{(1-x)^2}{1+x^2}.
\end{equation}
In particular, the exponent of $|x_i|$ in either majorant is exactly
\[
 (C-1)+g-(n-1)=C+g-n=a+2g;
\]
there is no singular finite-size correction at $x=0$.
Keeping the other powers as well gives the following exact identity:
\begin{align}
 \frac{\log\mathcal I_{\kappa,n}}{n^2}-\frac12\log2
 &=c_{\kappa,n}\log2
   +\left\langle W_\kappa(x)+\frac{\kappa+2}{2n}\log(1+x^2)\right\rangle
   +\langle W_s(s)\rangle\notag\\
 &\quad+\frac\kappa2\langle\log|x-x'|\rangle_{\ne}
       +\frac12\langle\log(1-xx')\rangle_{\ne}
       +\langle\log|s-s'|\rangle_{\ne}\notag\\
 &\quad-\langle\log(1-2xs+x^2)\rangle,
 \label{eq:energy-exact-finite}\\
 c_{\kappa,n}&=\frac Cn+\frac{\kappa-2}{2}
                       -\frac{\kappa+1}{2n}.\notag
\end{align}
The displayed correction involving $\log(1+x^2)$ is bounded by
$2\log2/n$ on the whole domain.

We use the classical Chebyshev expansion of the logarithmic kernel,
in the form of Haagerup's identity presented in
\cite[Lemma~3.1 and its proof]{GaroufalidisPopescu2013}.
We first record its damped form on $[-1,1]$. If
$u=\cos\theta$, $u'=\cos\theta'$, and $0\le r<1$, set
\begin{align}
 L_r(u,u')
 &=-\log2+\log|1-re^{i(\theta+\theta')}|
                 +\log|1-re^{i(\theta-\theta')}|\notag\\
 &=-\log2-2\sum_{k\ge1}\frac{r^kT_k(u)T_k(u')}{k}.
 \label{eq:energy-cosine-kernel}
\end{align}
This follows from the power series for $\log(1-z)$ and the addition
formula for cosines.  Letting $r\uparrow1$ at distinct $u,u'$ gives
$L_1(u,u')=\log|u-u'|$, since the product of the two chord lengths
is $2|u-u'|$.  The other two identities, for $|z|,|z'|<1$, are
\begin{equation}\label{eq:energy-power-cross-kernels}
 \log(1-zz')=-\sum_{k\ge1}\frac{z^kz'^k}{k},\qquad
 -\log(1-2zs+z^2)=2\sum_{k\ge1}\frac{z^kT_k(s)}{k}.
\end{equation}
For the last identity the logarithm is the logarithm of a positive real
number, equal to $\log|1-ze^{i\arccos s}|^2$.
The singular kernels cannot be summed at empirical diagonals. We
regularize all appearances of the same moment by the same factor so
that the pure and cross interactions still complete a square. Near
$(x,s)=(1,1)$, the identity
$|1-xe^{i\arccos s}|^2=(1-x)^2+2x(1-s)$ identifies the scales
$1-x$ and $\sqrt{1-s}$ used in the following damping factors.

Fix $0<\varepsilon<1/8$ and define
\begin{equation}\label{eq:energy-damping}
 \tau=1-\varepsilon,\qquad
 \rho(x)=1-\varepsilon(1-x),\qquad
 \sigma(s)=1-\varepsilon\sqrt{1-s}.
\end{equation}
Replace the $x$ cosine kernel in Equation~\eqref{eq:energy-exact-finite}
by $L_{\tau^2}(x,x')$ and the $s$ cosine kernel by
$L_{\sigma(s)\sigma(s')}(s,s')$.  Replace the power kernel by
\[
 \log(1-\rho(x)\rho(x')xx'),
\]
and the cross kernel by
\[
 -\log|1-\rho(x)\sigma(s)xe^{i\arccos s}|^2.
\]
Every original kernel is bounded above by its replacement plus
$O(\varepsilon)$, with an absolute constant independent of the configuration.
We give the global estimates, including the signs in the power kernel.

For $|\zeta|=1$ and $0<r\le1$,
\[
 |1-r\zeta|^2-r|1-\zeta|^2=(1-r)^2\ge0.
\]
Thus each of the cosine kernels costs at most $-\log r$ in the
comparison, and here $r\ge(1-\varepsilon)^2$.  For the power kernel put
$u=xx'$ and $r=\rho(x)\rho(x')$.  If $u\ge0$, monotonicity gives
$\log(1-u)\le\log(1-ru)$.  If $u<0$, write $q=-u\le1$; then
\begin{equation}\label{eq:energy-opposite-sign}
 \log(1+q)-\log(1+rq)
 \le\frac{(1-r)q}{1+rq}\le1-r\le4\varepsilon.
\end{equation}
For the cross kernel write
$a=1-x$, $b=\sqrt{1-s}$, $\theta=\arccos s$,
$z=1-xe^{i\theta}$, and $z_\varepsilon=1-\rho(x)\sigma(s)xe^{i\theta}$.
Then
$|z_\varepsilon-z|\le\varepsilon|x|(a+b)$.
On $x\le0$ we have $|z|\ge1$ and $|x|(a+b)\le3$.
On $0\le x\le1/2$ we have $|z|\ge1/2$ and
$|x|(a+b)\le1$.  On $1/2\le x<1$ we have
\[
 |z|^2=a^2+2xb^2\ge a^2+b^2,
 \qquad |x|(a+b)\le\sqrt2\,|z|.
\]
Consequently $|z_\varepsilon-z|\le3\varepsilon|z|$ on all three domains,
and in the direction required for an upper bound,
\begin{equation}\label{eq:energy-cross-comparison}
 -\log|z|^2\le-\log|z_\varepsilon|^2
                +2\log(1+3\varepsilon)
 \le-\log|z_\varepsilon|^2+6\varepsilon.
\end{equation}

For each fixed finite interior configuration all the damped series are
absolutely convergent.  Put
\[
 A_k=\langle\tau^kT_k(x)\rangle,\qquad
 B_k=\langle\rho(x)^kx^k\rangle,\qquad
 C_k=\langle\sigma(s)^kT_k(s)\rangle.
\]
Inserting the diagonals gives the two exact negative-square identities
\begin{align}
 \frac\kappa2\langle L_{\tau^2}(x,x')\rangle
 &=-\frac\kappa2\log2-\kappa\sum_{k\ge1}\frac{A_k^2}{k},
 \label{eq:energy-first-square}
\end{align}
\begin{align}
 &\frac12\langle\log(1-\rho(x)\rho(x')xx')\rangle
 +\langle L_{\sigma(s)\sigma(s')}(s,s')\rangle\notag\\
 &\qquad-\langle\log|1-\rho(x)\sigma(s)xe^{i\arccos s}|^2\rangle\notag\\
 &\quad=-\log2-\frac12\sum_{k\ge1}\frac{(B_k-2C_k)^2}{k}.
 \label{eq:energy-second-square}
\end{align}
In particular the $B_k$ and $C_k$ in the cross term are exactly the same
ones as in the pure power and cosine terms.  Adding the constants in
these identities to $c_{\kappa,n}\log2$ gives
\[
 \left(-1+\alpha+\gamma-\frac{\kappa+1}{2n}\right)\log2.
\]
The omitted-diagonal constants are part of the next $O_\varepsilon(1/n)$
correction.

We control this correction before taking any supremum.  A damped cosine
diagonal is bounded below by $-\log2+2\log(1-r)$, so the $x$ cosine
diagonals cost $O_\varepsilon(1/n)$.  For the power diagonal,
\[
 1-\rho(x)^2x^2\ge1-x\quad(x\ge0),\qquad
 1-\rho(x)^2x^2\ge1-(1-\varepsilon)^2\ge\varepsilon\quad(x\le0).
\]
Since $1\le1-x\le2$ when $x\le0$, their upward correction is at most
\[
 O_\varepsilon(1/n)-\frac1{2n}\langle\log(1-x)\rangle.
\]
Finally,
\[
 1-\sigma(s)^2
 =\varepsilon\sqrt{1-s}\,(2-\varepsilon\sqrt{1-s})
 \ge\varepsilon\sqrt{1-s},
\]
so the $s$ cosine diagonals cost at most
$O_\varepsilon(1/n)-n^{-1}\langle\log(1-s)\rangle$.
The cross term already contains all pairs.  There is no loss involving
$\log|x|$ or $\log(1+x)$.

For real numbers $a,b$ we have $-a^2\le b^2-2ab$.  Apply this to the
first square with $b=p_k$, and to the second with $b=v_k$; after dividing
by $k$ and summing, the two inequalities are
\begin{align*}
 -\kappa\sum_k A_k^2/k
 &\le\kappa\|p\|_*^2-2\kappa\sum_kp_kA_k/k,\\
 -\tfrac12\sum_k(B_k-2C_k)^2/k
 &\le\tfrac12\|v\|_*^2-\sum_kv_kB_k/k+2\sum_kv_kC_k/k.
\end{align*}
For $\kappa=1$, failure of Equation~\eqref{eq:case-condition} says
\[
 1-2\left\langle\frac{x^2}{1+x^2}\right\rangle
 >1-\frac1n-2\gamma,
 \qquad\text{hence}\qquad \langle D(x)\rangle>-1/n.
\]
It follows that adding $\lambda\langle D(x)\rangle$ costs at most
$\lambda/n$ for an upper bound; this explains both $\lambda\ge0$ and its
sign in Equation~\eqref{eq:energy-dual}.

Let $T_\varepsilon(p,x)$, $S_\varepsilon(v,x)$, and
$T_\varepsilon(v,s)$ denote the three series in
Equation~\eqref{eq:energy-trial-notation} with factors
$\tau^k$, $\rho(x)^k$, and $\sigma(s)^k$, respectively.
The preceding pointwise bounds reduce the upper estimate to the sum of
two one-variable suprema, with functions
\begin{align*}
 X_{n,\varepsilon}(x)
 &=\frac{19}{48}\log|x|
   +\left(\frac1{12}-\frac1{2n}\right)\log(1-x)
   \\
 &\quad-(\kappa/2+\alpha+\eta+\gamma)\log(1+x^2)
   +\lambda D(x)\\
 &\quad-2\kappa T_\varepsilon(p,x)-S_\varepsilon(v,x),\\
 Y_{n,\varepsilon}(s)
 &=\frac7{48}\log s
   +\left(\frac1{12}-\frac1n\right)\log(1-s)
   +2T_\varepsilon(v,s),
\end{align*}
plus the constant in Equation~\eqref{eq:energy-dual} and
$O(\varepsilon)+O_\varepsilon(1/n)$.  All bounded finite-size terms from
Equation~\eqref{eq:energy-exact-finite} are included in this last error.
For $n\ge24$ both weakened endpoint coefficients are at least $1/24$.

Set $P_1=\sum_k|p_k|$ and $V_1=\sum_k|v_k|$, which are finite.
The inequality $1-(1-d)^k\le kd$ shows uniformly that
\begin{align*}
 |T_\varepsilon(p,x)-T(p,x)|&\le\varepsilon P_1,\\
 |S_\varepsilon(v,x)-S(v,x)|&\le2\varepsilon V_1,\\
 |T_\varepsilon(v,s)-T(v,s)|&\le\varepsilon V_1.
\end{align*}
Thus the changes in the two tangent functions are at most
$2\kappa\varepsilon P_1+2\varepsilon V_1$ and
$2\varepsilon V_1$, respectively.  These are common summable bounds;
no uniform convergence of the raw square series is required.
All remaining nonsingular terms are bounded uniformly in $n\ge24$ and
$0<\varepsilon<1/8$.  The unchanged positive coefficients $19/48,7/48$
and the weakened coefficients at least $1/24$ force uniform decay to
$-\infty$ near $x=0$, $x=1$, $s=0$, and $s=1$.
The values at the fixed comparison points $x=-1/2$, $s=1/2$ have a
common finite lower bound. For the first supremum we include the
regular endpoint $x=-1$ by continuous extension.
There are therefore compact sets, independent
of $n\ge24$ and small $\varepsilon$, containing maximizers for both
suprema. On these sets the functions converge uniformly as $n\to\infty$
with $\varepsilon$ fixed, and subsequently as $\varepsilon\downarrow0$.
Taking the limits in precisely this order removes
$O_\varepsilon(1/n)$ first and proves Equation~\eqref{eq:energy-dual}.
Finally apply Equation~\eqref{eq:real-supremum-reduction}.
\end{proof}
